\begin{frame}{One set of settings}
\vspace{.8em}
\centering
{\Large Integral threshold: \textbf{\textcolor{pennred}{0.3}}}\par
\vspace{1.8em}
{\Large Vertex finding: \textbf{\textcolor{pennred}{$|\Delta z|\leq0.1$ mm}}}\par
\vspace{1.8em}
{\Large Vertex categories: \textbf{\textcolor{pennred}{ATLAS convention}}}\par
\vspace{.45em}
{\small Clean / Matched, Merged, Split, Fake}\par
\vfill
\sourcecaption{0.5 is a cross-check only. Current category counts approximate ATLAS: unit track weights and an unlabelled-track proxy for Fake.}
\end{frame}

\begin{frame}{HL-LHC: resolution at integral cuts 0.3 and 0.5}
\centering
\includegraphics[width=.99\linewidth,height=.78\textheight,keepaspectratio]{figures/resolution_integral_030_050.pdf}
\sourcecaption{AMVF is unchanged between panels. This reconstructed-pair fit uses no truth-matching window.}
\end{frame}

\begin{frame}{One vertex classification: the ATLAS convention}
\context{Same four categories as V4 slide 26; truth contributions use vertex-fit weights.}
\small\renewcommand{\arraystretch}{1.5}
\begin{tabular}{p{3.2cm}p{10.3cm}}
\toprule
\textbf{Clean / Matched} & At least 70\% of the total track weight comes from one true collision.\\
\textbf{Merged} & No true collision contributes 70\% of the total track weight.\\
\textbf{Split} & Several reco vertices have the same dominant true collision. Only the vertex with the largest $\sum p_T^2$ keeps Clean or Merged.\\
\textbf{Fake} & Fake tracks contribute more weight than any single true collision.\\
\bottomrule
\end{tabular}
\par\vspace{.6em}\footnotesize
Fake and Split take precedence over the Clean/Merged label. We retain reco vertices with at least two tracks for both methods.
\sourcecaption{\amvfnote. Current plots use unit weights and unlabelled tracks as a Fake proxy; the ntuples lack fit weights and detailed fake-track truth information. Position-only results use ``nearby'' and ``unmatched'', not these category names.}
\end{frame}
